% 题干转写：完成
% 解析转写：完成
% 图形核对：完成
\section{第37届复赛第6题：光纤陀螺仪与Sagnac效应}

\subsection{题干}
光纤陀螺仪利用Sagnac效应测量转动角速度。半径为 $R$ 的光纤环路以角速度 $\Omega$ 逆时针转动，固定在环上的分束器 $A$ 将相干光分成沿逆时针（$\mathrm{CCW}$）和顺时针（$\mathrm{CW}$）传播的两束光，如图\ref{fig:cpho37q06a}所示。真空光速为 $c$。
\begin{figure}[htbp]\centering
\begin{tikzpicture}[scale=0.86,>=Stealth,font=\small]
  \draw[thick] (0,0) circle (2.05);
  \fill (0,0) circle (1.5pt);
  \draw[->,thick] (0,0) -- (2.05,0) node[midway,below] {$R$};
  \node[above] at (0,2.05) {$A$};
  \draw[thick] (0,2.05) -- (0,1.62);
  \draw[->,thick] (-0.62,1.95) arc[start angle=108,end angle=145,radius=2.05];
  \node[above left] at (-1.50,1.52) {$\mathrm{CCW}$};
  \draw[->,thick] (0.62,1.95) arc[start angle=72,end angle=35,radius=2.05];
  \node[above right] at (1.50,1.52) {$\mathrm{CW}$};
  \draw[->,thick] (-0.38,-0.30) arc[start angle=230,end angle=112,radius=0.56];
  \node[right] at (0.18,0.37) {$\Omega$};
\end{tikzpicture}

\caption{第37届复赛第6题光纤环路、分束器与反向传播光的TikZ重绘。}\label{fig:cpho37q06a}
\end{figure}

（1）如图\ref{fig:cpho37q06b}，光纤由折射率分别为 $n_1$、$n_2$ 的内、外两层介质构成，且 $n_1>n_2$。端口外介质折射率为 $n_0$。求使光能在光纤内传输时，输入端入射角 $i$ 的条件。
\begin{figure}[htbp]\centering
\begin{tikzpicture}[x=0.85cm,y=0.85cm,>=Stealth,font=\small]
  \draw[thick] (-3.4,0) -- (4.7,0);
  \draw[densely dashed] (-0.9,0.82) -- (3.45,0.82);
  \draw[densely dashed] (-0.9,-0.82) -- (3.45,-0.82);
  \draw[thick] (-1.15,0) ellipse (0.40 and 0.92);
  \draw[thick] (-0.78,0) ellipse (0.30 and 0.62);
  \draw[thick] (3.72,0) ellipse (0.40 and 0.92);
  \draw[densely dashed] (3.35,0) ellipse (0.30 and 0.62);
  \fill (-1.15,0) circle (1.5pt);
  \fill (3.72,0) circle (1.5pt);
  \draw[->,thick] (-3.05,-1.02) -- (-1.15,0);
  \draw (-2.55,0) arc[start angle=0,end angle=28,radius=0.60];
  \node[below] at (-2.35,-0.18) {$i$};
  \node[above] at (-2.65,0.16) {$n_0$};
  \node[above] at (2.75,0.82) {$n_1$};
  \node[below] at (2.75,-0.82) {$n_2$};
\end{tikzpicture}

\caption{第37届复赛第6题双层光纤输入端结构的TikZ重绘。}\label{fig:cpho37q06b}
\end{figure}

（2）沿环路密绕 $N$ 匝光纤，首尾相接于分束器 $A$。介质折射率为 $n_1$，真空波长为 $\lambda_0$。求两束光的传播时间 $t_{\mathrm{CCW}}$、$t_{\mathrm{CW}}$ 和相位差，并说明相位差与 $n_1$ 的关系。

（3）谐振式光学陀螺仪含有沿逆时针方向转动的光学环形腔，半径为 $R$，角速度为 $\Omega$，且 $R\Omega\ll c$。腔中介质折射率为 $n_1$。静止时同一级次 $m$ 的模式满足 $2\pi R=m\lambda_{m0}$。求旋转时顺、逆时针共振模式频率差 $\Delta f$ 与 $\Omega$ 的关系。

\subsection{参考解答整理}
\paragraph{（1）光纤端口入射条件。}
设光在端口处的折射角为 $i'$，到内外层界面时的入射角为 $i''$。折射定律、全反射条件和几何关系分别为
\[
n_0\sin i=n_1\sin i',\qquad
i''\geqslant i_c,\qquad \sin i_c=\frac{n_2}{n_1},\qquad
i'+i''=\frac\pi2.
\]
故
\[
n_0\sin i=n_1\cos i''\leqslant n_1\cos i_c
=\sqrt{n_1^2-n_2^2},
\]
即
\[
\boxed{i\leqslant\arcsin\frac{\sqrt{n_1^2-n_2^2}}{n_0}.}
\]

\paragraph{（2）Sagnac相位差。}
在实验室系中，两束光的线速度为
\[
v_{\mathrm{CCW}}=
\frac{c/n_1+\Omega R}{1+(c/n_1)(\Omega R)/c^2},
\qquad
v_{\mathrm{CW}}=
\frac{c/n_1-\Omega R}{1-(c/n_1)(\Omega R)/c^2}.
\]
环路随分束器运动，故传播时间满足
\[
v_{\mathrm{CCW}}t_{\mathrm{CCW}}=N2\pi R+\Omega Rt_{\mathrm{CCW}},
\qquad
v_{\mathrm{CW}}t_{\mathrm{CW}}=N2\pi R-\Omega Rt_{\mathrm{CW}}.
\]
图\ref{fig:cpho37q06sagnac}标出了传播期间分束器由 $A(t=0)$ 移至 $A(t)$ 的弧长 $\Omega Rt$，这正是两个传播路程出现正、负修正项的几何原因。
\begin{figure}[htbp]\centering
\begin{tikzpicture}[x=.72cm,y=.72cm,>=Stealth,font=\scriptsize]
  \draw[thick] (0,0) circle (2.05);
  \coordinate (A) at (0,-2.05);
  \coordinate (At) at (.95,-1.82);
  \fill (A) circle(1.5pt) node[below]{$A(t=0)$};
  \fill (At) circle(1.5pt) node[right]{$A(t)$};
  \draw[->,thick] (A) arc(-90:42:2.05) node[midway,right]{$\Omega Rt$};
  \draw[->,very thick] (-1.45,-1.45) arc(-135:70:2.05) node[midway,left]{$\mathrm{CCW}$};
  \draw[->,very thick] (1.45,-1.45) arc(-45:-250:2.05) node[midway,right]{$\mathrm{CW}$};
  \draw[->] (2.6,.15) arc(0:75:.65) node[midway,right]{$\Omega$};
  \node at (0,0) {半径 $R$};
  \draw[densely dashed] (A)--(At);
\end{tikzpicture}

\caption{第37届复赛第6题Sagnac效应中分束器随环路转动的路径补偿示意图。}\label{fig:cpho37q06sagnac}
\end{figure}
解得
\[
\boxed{t_{\mathrm{CCW}}=
\frac{N2\pi R(n_1c+\Omega R)}{c^2-(\Omega R)^2},\qquad
t_{\mathrm{CW}}=
\frac{N2\pi R(n_1c-\Omega R)}{c^2-(\Omega R)^2}.}
\]
于是相位差为
\[
\boxed{\Delta\varphi=
\frac{2\pi c}{\lambda_0}(t_{\mathrm{CCW}}-t_{\mathrm{CW}})
=\frac{8\pi^2NR^2\Omega c}{\lambda_0[c^2-(\Omega R)^2]}.}
\]
由此可见，在本题模型中 $\Delta\varphi$ 与光纤介质折射率 $n_1$ 无关。

\paragraph{（3）谐振模式频率分裂。}
对同一级次 $m$ 的两类共振模式，分别有
\[
L_{\mathrm{CW}}=m\lambda_{\mathrm{CW}},\qquad
L_{\mathrm{CCW}}=m\lambda_{\mathrm{CCW}}.
\]
由旋转腔中的传播路程可得频率差
\[
\Delta f=f_{\mathrm{CW}}-f_{\mathrm{CCW}}
=\frac{m\Omega}{\pi}
\frac{c^2-(\Omega R)^2}{c^2-(n_1\Omega R)^2}.
\]
在 $R\Omega\ll c$ 下，忽略二阶小量，
\[
\Delta f\approx\frac{m\Omega}{\pi}.
\]
故环形腔转动角速度可写为
\[
\boxed{\Omega=\frac\pi m\Delta f
=\frac{\lambda_{m0}}{2R}\Delta f.}
\]

\subsection{转写依据}
题干、图6a、图6b及公式均依据本题题面转写稿、答案公式核对表和图件核对入口。两幅图均以TikZ重绘，不使用原题截图。本题第（1）问10分，第（2）问18分，第（3）问12分。
